\section{Honest Gaps}
\label{sec:gaps}

GPU-native verify is the LP-203 architecture. The matrix bench on
2026-06-04 demonstrated that the architecture is not yet
end-to-end exercised on the production hot path. This section
inventories exactly what is not measured, why, and what closes
each gap.

\subsection{The matrix bench's GPU-utilization observation}

The Quasar 4-scheme matrix bench \cite{matrixbench} ran for two
hours on Spark on 2026-06-04 and captured \num{4595} samples of
\texttt{nvidia-smi dmon} during the run. The aggregate
observation: mean \texttt{SM\%} = \SI{11.9}{\percent}.

\textbf{The GPU was not exercised by any path in this bench.}
Every cert mode (PQ-off, PQ-fast, PQ-strict, PQ-heavy) and every
$N{=}\{4,8,16,32,64\}$ ran on CPU. The bench's measured
wall-clock numbers (PQ-strict at \SI{595}{\milli\second}
$N{=}64$, see LP-217 \cite{lp217}) are CPU numbers. The GPU was
idle (or running background work) for the entire run.

This is not a benchmarking error. It is an honest reflection of
where the production code stops dispatching to GPU today.

\subsection{Boundary gap A: \texttt{lux-lattice.pc} missing on Spark}

The first boundary gap is at the cgo / pkg-config layer. The
CUDA target's dependency chain requires \texttt{lux-lattice.pc}
to be findable by \texttt{find\_package(MLX REQUIRED)} in the
CMake build of \texttt{lattice/v7/gpu}. On Spark today,
\texttt{lux-lattice.pc} is not installed; the
\texttt{find\_package(MLX REQUIRED)} call fails; the
\texttt{lattice/v7/gpu} target does not build under
\texttt{-tags "cgo gpu"}.

\paragraph{Consequence.} Even if a caller-side path tried to
dispatch lattice ops to the GPU, the CUDA backend's lattice path
is missing from the binary. The caller falls through to the CPU
fallback. The mean \texttt{SM\%} of \SI{11.9}{\percent} is the
direct expected consequence.

\paragraph{Close-out.} Install \texttt{lux-lattice.pc} on Spark
(packaged from \texttt{lux-private/lattice/v7/pkgconfig}). Once
the package is present, \texttt{find\_package(MLX REQUIRED)}
resolves and the lattice GPU target builds. No code change
required on either side.

\subsection{Boundary gap B: Magnetar \texttt{lux\_slhdsa\_*}
entry points not exposed}

The second boundary gap is at the C API surface. The Magnetar
SLH-DSA \cite{lp181,fips205} CUDA kernel exists on Spark and
passes the FIPS 205 KAT vectors when called directly. But the
\texttt{lux\_slhdsa\_*} entry points are not exposed in
\texttt{lux/accel/c\_api.h}. \texttt{luxfi/accel} has no symbol
to dispatch to; the dispatcher falls back to CPU.

\paragraph{Consequence.} Magnetar verify in the PQ-heavy and
strict-PQ-heavy modes runs on CPU even when the GPU has a
working kernel. SLH-DSA-192s per-validator sign on Spark CPU is
$\sim$\SI{795}{\milli\second}; the GPU kernel would be
substantially faster if dispatched, but the C API surface
prevents dispatch.

\paragraph{Close-out.} Add \texttt{lux\_slhdsa\_keygen},
\texttt{lux\_slhdsa\_sign}, \texttt{lux\_slhdsa\_verify} to
\texttt{c\_api.h}; wire them into
\texttt{luxfi/accel/ops\_lattice.go} \cite{accelmod}. The kernels
on the other side of the ABI are already there.

\subsection{Composite signature-verify wrappers as stubs}

The CUDA plugin in
\texttt{lux-private/gpu-kernels/backends/cuda/src/plugin.cpp}
exposes 76 vtbl ops. The honest split:

\paragraph{Real GPU code with measured throughput.}
\begin{itemize}[leftmargin=2em,nosep]
\item BLS12-381 EIP-2537: G1 add, G1 MSM, G2 add, G2 MSM,
  pairing, hash-to-curve
\item BN254: G1 ops, pairing-check
\item keccak256, sha3\_256, blake3
\item modexp
\item Magnetar SLH-DSA primitives: \texttt{wotsplus\_chain},
  \texttt{fors\_subtree}, \texttt{xmss\_subtree},
  \texttt{hmsg\_prfmsg}
\item Corona partial-sign step
\end{itemize}

\paragraph{Explicit \texttt{LUX\_BACKEND\_ERROR\_NOT\_SUPPORTED}
stubs (16 of 76 ops).}
\begin{itemize}[leftmargin=2em,nosep]
\item Direct-verify entries for signature schemes: Ed25519,
  sr25519, ML-DSA, SLH-DSA composite, ML-KEM
\item Projective-coords BLS / BN254 ops (the EIP-2537 path is
  the canonical surface; projective coords are absent)
\end{itemize}

\paragraph{The composite signature-verify wrappers are the gap
to close before LP-203 numbers turn fully empirical.} The
primitives (pairings, NTT, hashing, Magnetar 4) are real and
measured; the wrappers that compose them into a single
``verify this signature'' call return \texttt{NOT\_SUPPORTED}.

\subsection{What the LP-203 markdown projects vs. what the
bench measured}

The LP-203 markdown's throughput tables include both
``measured'' and ``estimated'' columns. The 2026-06-04 bench
forces a reassessment of one specific number:

\begin{table}[h]
\centering
\small
\begin{tabular}{@{}p{4.4cm}rrp{2.6cm}@{}}
\toprule
\textbf{Operation} & \textbf{LP-203 est.}
  & \textbf{GB10 meas.} & \textbf{Delta} \\
\midrule
BLS verify, $N{=}1$, H100 FP64-roofline
  estimate & \SI{0.5}{\micro\second}
  & \SI{19}{\micro\second}
  & $\sim 38\times$ slower than roofline \\
BLS verify, 1000-sig batch, H100 estimate
  & \SI{0.5}{\milli\second}
  & $\sim$\SI{20}{\milli\second}
  & $\sim 40\times$ slower than roofline \\
\bottomrule
\end{tabular}
\caption{Measured GB10 BLS pairing cost is approximately
$40\times$ the LP-203 H100 FP64 roofline projection. The
measurement is the Miller-loop-plus-final-exp cost on real
sm\_120 silicon, not arithmetic intensity over peak FLOPS. The
GB10 number is the correct figure for downstream LP
performance budgets; the H100 roofline projection should be
re-derived against the measured GB10 number, not used directly.}
\label{tab:projection-delta}
\end{table}

\subsection{What is NOT yet measured}

\begin{itemize}[leftmargin=2em,nosep]
\item \textbf{Pulsar / Magnetar composite verify GPU
  throughput.} The kernels exist; the composite-verify wrappers
  are stubs. Until the wrappers ship, the production hot path
  for ML-DSA and SLH-DSA verify runs on CPU. The matrix bench's
  mean \texttt{SM\%} of \SI{11.9}{\percent} is the empirical
  consequence.
\item \textbf{GPU dispatch on Spark blocked by
  \texttt{lux-lattice.pc} absence.} Until the package is
  installed, no lattice GPU code is in the binary; the CPU
  fallback path is the only path.
\item \textbf{CXL 3.0 racks.} The LP-203 markdown does not
  reference CXL 3.0, but a sibling LP-205 will -- the bench
  cluster does not yet include CXL 3.0 hosts, so cross-socket
  UMA semantics under CXL are not measured.
\item \textbf{End-to-end block finality measurements on the
  full LP-200 stack with GPU dispatch.} The matrix bench
  measured Quasar at the threshold-cert layer; full pipeline
  numbers (NIC ingest $\to$ verify $\to$ Merkle $\to$ ZapDB
  commit) await the wrapper close-out.
\end{itemize}

\subsection{What this means for downstream LPs}

LP-217 \cite{lp217} (cert profile modes), LP-218 \cite{lp218}
(P3Q rollups), LP-205 (pipelined Quasar): every paper that
cites a GPU latency for ML-DSA or SLH-DSA verify is citing a
\emph{projection}, not a measurement. The papers tag those
cells ``projected, not measured'' for that reason. LP-217's
matrix-bench measurements are CPU wall-clock at $N{=}64$ on
Spark, not GPU. LP-218's P3Q verify cost of
\SI{50}{\micro\second} on Blackwell is a projection.

This is not a problem with the LPs. The LPs are honest about
what is measured and what is projected. It is a problem with
the dispatch surface, which two small CL changes
(\texttt{lux-lattice.pc} install plus
\texttt{lux\_slhdsa\_*} C API exposure) close.

\subsection{Pulsar's per-validator sign measurement}

One CPU measurement worth recording, because it sets the
GPU-target ceiling: Pulsar per-validator sign on Spark CPU is
\SI{175}{\micro\second} \cite{paritytest}. The GPU target for
the Pulsar sign kernel (when it lands behind a working
composite-verify wrapper) should beat that number. Pulsar's
algorithm \cite{lp073} is Shamir-on-GF($257$) plus one FIPS 204
sign \cite{fips204} on the aggregator; the FIPS 204 sign is the
work to GPU-accelerate.
